\section{General measures in the minimal class}
\label{sec:big-measure-smoothing}

We now solve the equations in $\mathcal E_V$. For a measure $\nu$ with $A_\theta(\nu)<\infty$, the coordinate convolution of \cref{subsec:psh-measure-smoothing} gives approximating equations with smooth densities. The integral estimates for this smoothing, combined with the higher-moment bound, give convergence of the solutions against the original measure. The resulting self-measure error yields convergence in capacity, and a first-moment tail estimate identifies the limit of the Monge--Amp\`ere measures.

\Needspace{18\baselineskip}
\begin{proposition}\label{prop:big-smoothed-measure-continuity}
Let $\nu$ be a non-pluripolar probability measure with $A_\theta(\nu)<\infty$, and let $\nu_j$ be its smooth approximations defined by \eqref{eq:measure-smoothing-duality}. Fix a constant $C>0$ and let $v_j,v\in\PSH(X,\theta)$ satisfy
\[
 v_j\sim V,\qquad \sup_Xv_j=0,\qquad
 \mu_{v_j}\leq C\nu_j,
 \qquad v_j\longrightarrow v\text{ in }L^1(X).
\]
Then $v\in\mathcal E_V$, the convergence also holds in $L^1(\nu)$ and in capacity, and
\[
 \mu_{v_j}\rightharpoonup\mu_v.
\]
\end{proposition}
\begin{proof}
The limit $v$ is not identically $-\infty$ and satisfies $\sup_Xv=0$, and $v_j,v\leq V\leq0$. By \cref{lem:measure-smoothing-integrals}, $\sup_jA_\theta(\nu_j)<\infty$, which gives the uniform estimate
\begin{equation}\label{eq:smoothed-first-moment}
 \int_X(-v_j)\,\mu_{v_j}
 \leq C\int_X(-v_j)\,\mathrm{d}\nu_j\leq C'.
\end{equation}
In particular the relative first moments $\int_X(V-v_j)\mu_{v_j}$ are bounded, and \cref{prop:big-relative-first-moment}, with fixed reference $w=V$, puts $v$ in $\mathcal E_V$.

By \cref{prop:big-minimal-higher-moment}, applied to the fixed measure $\nu$, the $v_j$ are uniformly integrable with respect to $\nu$. For each fixed $T$, the bounded functions $\max(v_j,-T)$ have a common lower curvature bound in a Hermitian background $\Omega\geq\theta$ and converge in $L^1$ for smooth volume to $\max(v,-T)$; hence \cref{lem:fixed-npp-integrals} gives their convergence in $L^1(\nu)$. These truncations need not be $\theta$-plurisubharmonic; they are used only in the larger Hermitian background. Uniform integrability, followed by $T\to\infty$, gives $v_j\to v$ in $L^1(\nu)$, and the limit is integrable because $A_\theta(\nu)<\infty$.

The smoothing lemma, in the form \eqref{eq:convolution-moving-integrals}, now yields
\[
 \int_X|v_j-v|\,\mu_{v_j}
 \leq C\int_X|v_j-v|\,\mathrm{d}\nu_j\longrightarrow0.
\]
Therefore the self-measure convergence theorem \cref{thm:big-relative-self-measure-convergence} gives convergence in capacity. Finally, \eqref{eq:smoothed-first-moment} gives the sublevel estimate
\[
 \sup_j\mu_{v_j}(\{v_j\leq-T\})\leq C'/T,
\]
and the masses of the $\mu_{v_j}$ are bounded by $M_V$. The tail criterion \cref{thm:general-background-damping} therefore gives the asserted weak convergence.
\end{proof}

\Needspace{36\baselineskip}
\subsection{Exponential and normalized equations}
\label{sec:big-minimal-general-measures}

The smoothing argument solves both equations for measures with finite $A_\theta$. A bounded carrier representation reduces an arbitrary finite non-pluripolar measure to increasing bounded densities, whose solutions have measures dominated by one fixed finite measure; \cref{thm:big-relative-fixed-continuity} then passes to the limit.

\begin{theorem}\label{thm:big-minimal-general-measures}
For every finite non-pluripolar measure $\mu$ with $\mu(X)>0$, the following assertions hold.
\begin{enumerate}
\item For every $\lambda>0$ there is a unique $u\in\mathcal E_V$ such that
\[
 \mathrm{MA}_\theta(u)=\mathrm{e}^{\lambda u}\mu.
\]
\item There are $v\in\mathcal E_V$ with $\sup_Xv=0$ and a unique scalar $c>0$ such that
\[
 \mathrm{MA}_\theta(v)=c\mu.
\]
\end{enumerate}
\end{theorem}
\begin{proof}
First suppose that $\mu$ is a probability measure with $A_\theta(\mu)<\infty$, and let $\mu_j$ be its smooth probability approximations. The theorems for the minimal ordinary type, \cref{thm:big-minimal-exponential,thm:big-minimal-normalized}, apply to their smooth, nonzero densities.

For the normalized equation, they give $v_j\sim V$ with $\sup_Xv_j=0$ and $\mu_{v_j}=c_j\mu_j$. Since $\mu_j$ is a probability measure, the mass bounds $\ell_V\leq\mu_{v_j}(X)\leq M_V$ give $\ell_V\leq c_j\leq M_V$. After extraction, $c_j\to c>0$ and $v_j\to v$ in $L^1(X)$. The preceding proposition applies, so $v\in\mathcal E_V$ and $\mu_{v_j}\rightharpoonup\mu_v$. The right-hand sides converge weakly to $c\mu$, which proves the equation.

For the positive-exponent equation, take $u_j\sim V$ satisfying $\mu_{u_j}=\mathrm{e}^{\lambda u_j}\mu_j$, and write $a_j=\sup_Xu_j$ and $v_j=u_j-a_j$. The mass bounds and Jensen's inequality give
\[
 \ell_V\leq \mathrm{e}^{\lambda a_j},\qquad
 M_V\geq \mathrm{e}^{\lambda a_j}\int_X\mathrm{e}^{\lambda v_j}\,\mathrm{d}\mu_j
       \geq \mathrm{e}^{\lambda a_j-\lambda A_\theta(\mu_j)},
\]
and $\sup_jA_\theta(\mu_j)<\infty$ by \cref{lem:measure-smoothing-integrals}. Thus $a_j$ is bounded above and below. After extraction, $a_j\to a$, $v_j\to v$ in $L^1(X)$, and $\mu_{v_j}\leq C\mu_j$. The preceding proposition gives $v\in\mathcal E_V$, convergence in $L^1(\mu)$ and weak convergence of the Monge--Amp\`ere measures. Set $u=v+a$. By the smoothing lemma, the right-hand sides converge weakly to $\mathrm{e}^{\lambda u}\mu$. This proves the equation, and $u\in\mathcal E_V$.

The same conclusions hold for every nonzero finite measure whose normalization to a probability measure has finite $A_\theta$, since scaling the measure only changes the positive scalar, or translates the solution of the exponential equation by a constant.

Now let $\mu$ be an arbitrary finite non-pluripolar measure of positive mass. The bounded-carrier decomposition \cref{lem:bounded-measure-carrier}, applied to $\omega$, writes
\[
 \mu=f\nu,\qquad \nu=\mathrm{MA}_\omega(\psi),\qquad
 \psi\in\PSH(X,\omega)\cap L^\infty(X),\qquad
 0\leq f\in L^1(\nu).
\]
For $\mu_k=\min(f,k)\nu$, we have $\mu_k(X)>0$. By \cref{lem:bounded-carrier-integrability}, the normalization of $\mu_k$ has finite $A_\theta$ for each fixed $k$.

For the normalized equation, the preceding case gives $v_k\in\mathcal E_V$ with $\sup_Xv_k=0$ and $\mu_{v_k}=c_k\mu_k$, where
\[
 \frac{\ell_V}{\mu(X)}\leq c_k\leq\frac{M_V}{\mu_1(X)}.
\]
After extraction, $c_k\to c>0$ and $v_k$ converges in $L^1$ to some $v\not\equiv-\infty$. All the measures $\mu_{v_k}$ are bounded by the fixed finite non-pluripolar measure $M_V\mu/\mu_1(X)$. The fixed-dominator theorem \cref{thm:big-relative-fixed-continuity} gives $v\in\mathcal E_V$ and weak convergence to $\mu_v$. Since $c_k\mu_k\to c\mu$ in total variation, $\mu_v=c\mu$.

For the positive-exponent equation, let $u_k\in\mathcal E_V$ solve $\mu_{u_k}=\mathrm{e}^{\lambda u_k}\mu_k$. The argument that follows the construction of the approximating solutions in the positive-exponent part of the proof of \cref{thm:general-measure-minimal} applies, with \cref{cor:big-relative-orders,thm:big-relative-fixed-continuity} and $\ell_V$ in place of \cref{cor:relative-extended-exponential-order,thm:relative-fixed-measure-continuity} and $v_-(\theta)$. It shows that the $u_k$ decrease to some $u\in\mathcal E_V$ with $\mu_u=\mathrm{e}^{\lambda u}\mu$.

Finally, \cref{cor:big-relative-orders} with reference $V$, applied in both directions, gives uniqueness of the exponential solution and of the positive scalar.
\end{proof}

\Needspace{23\baselineskip}
Positive lower mass of a Hermitian form gives positive lower volume to the minimal reference of every smooth big background.

\begin{corollary}\label{cor:metric-pvp-minimal-volume}
Suppose that $v_-(\omega)>0$. Then $\underline{\operatorname{vol}}_\theta(V_\theta)>0$ for every smooth big real $(1,1)$-form $\theta$ on $X$; hence the preceding theorem applies to every such $\theta$.
\end{corollary}
\begin{proof}
Choose a strict barrier $\rho\in\PSH(X,\theta)$, so that $\theta_\rho\geq a\omega$ for some $a>0$. Then $\rho\in\PSH(X,\theta-a\omega)$. Volume addition (\cref{prop:volume-addition}), followed by monotonicity in the singularities (\cref{thm:volume-monotonicity}), gives
\[
 0<a^n v_-(\omega)
 =\underline{\operatorname{vol}}_{a\omega}(0)
 \leq\underline{\operatorname{vol}}_\theta(\rho)
 \leq\underline{\operatorname{vol}}_\theta(V_\theta).
\]
In the addition step the big background is $a\omega$, and the added positive current is $(\theta-a\omega)_\rho$. Bounded mass gives finiteness of the upper volume.
\end{proof}
